% Section 14: 02-examples | Part: 2 | Depends on: see docs/REWRITE_PLAN.md | Status: todo | Original lines: 2336-2365
%!TEX root = ../../../main.tex
\begin{example}[Vector calculus in $\GA(\R^3)$]
Let $f$ be a scalar field and $\mathbf a$ a vector field. Then
\[
 \vd f=\operatorname{grad}f,\qquad
 \vd\mathbf a=\vd\cdot\mathbf a+\vd\wedge\mathbf a=\operatorname{div}\mathbf a+\pss\,\operatorname{curl}\mathbf a ,
\]
using Proposition~\ref{prop:cross} ($\vd\wedge\mathbf a=\pss\,(\vd\times\mathbf a)$). Since $\operatorname{curl}\mathbf a=\dual{(\vd\wedge\mathbf a)}$ (Proposition~\ref{prop:cross}), the two classical identities are both instances of $\vd\wedge\vd\wedge{}=0$: $\operatorname{curl}\operatorname{grad}f=\dual{(\vd\wedge\vd f)}=0$, and, by \eqref{eq:dualduality},
$\operatorname{div}\operatorname{curl}\mathbf a=\vd\lc\dual{(\vd\wedge\mathbf a)}=\dual{(\vd\wedge\vd\wedge\mathbf a)}=0$. Finally $\vd^2=\vd\cdot\vd$ is the Laplacian.
\end{example}

\begin{example}[Maxwell's equations in one line]
In the spacetime algebra $\GA(\R^{1,3})$ the electromagnetic field is the bivector $F=\mathbf E+\pss\mathbf B$ (with $\mathbf E,\mathbf B$ the electric and magnetic fields written as relative vectors, i.e.\ as bivectors of spacetime) and (in suitable units) the four Maxwell equations are the single equation
\[
 \vd F=J ,\qquad\text{i.e.}\qquad \underbrace{\vd\wedge F=0}_{\dd F=0\ (\text{homogeneous})},\quad
 \underbrace{\vd\lc F=J}_{\text{inhomogeneous, source }J}.
\]
The two pieces are precisely the exterior derivative and the codifferential (a dual exterior derivative) of $F$. The first equation has a potential, $F=\vd\wedge A$, by the Poincar\'e lemma.
\end{example}

\begin{intuition}
In geometric algebra the derivative is \emph{literally a vector}, $\vd=e^i\partial_i$, so it multiplies a field $F$ with the geometric product, and the geometric product $ab=a\cdot b+a\wedge b$ splits into a lower-grade and a higher-grade part:
\[
 \underbrace{\vd F}_{\text{one derivative}}=\underbrace{\vd\lc F}_{\text{divergence-like (grade}\downarrow)}+\underbrace{\vd\wedge F}_{\text{curl-like = exterior derivative (grade}\uparrow)} .
\]
The exterior derivative is thus the ``wedge half'' of a single, invertible derivative. Duality then converts the wedge half into the contraction half, so that curl and divergence are \emph{the same operation seen through the dual} ($\hodge$): that is why Maxwell's equations, which look like two different kinds of equation in vector calculus, become one equation. Everything in this section is just $ab=a\cdot b+a\wedge b$ with $a$ replaced by $\vd$.
\end{intuition}
